\section{שיעור 14 — תורת ההפרעות הבלתי-תלויה בזמן (28.12.2025)}
% מבוסס על תוכן הרצאה – ע"פ קלט המקור
% נתמך ומורחב באופן פדגוגי

\subsection*{מטרות השיעור}
\begin{itemize}
\item להבין את הרעיון הבסיסי של תורת ההפרעות ומתי ניתן להשתמש בה.
\item לגזור את נוסחאות התיקון לאנרגיות ולמצבים עצמיים עד לסדר שני.
\item ללמוד לזהות תנאי תקפות ומקרים בהם התורה נכשלת.
\item ליישם את השיטה על דוגמאות קונקרטיות.
\end{itemize}

\subsection*{רקע נדרש}
\begin{itemize}
\item משוואת שרדינגר עצמאית בזמן: $\hat{H}\ket{\psi} = E\ket{\psi}$
\item סימון דיראק: $\ket{\psi}$, $\bra{\phi}$, $\braket{\phi|\psi}$, $\bra{\phi}\hat{A}\ket{\psi}$
\item שלמות של בסיס אורתונורמלי: $\sum_n \ket{n}\bra{n} = \hat{1}$
\item סדרות טיילור וכיצד להרחיב פונקציות בפרמטר קטן.
\end{itemize}

\subsection{מוטיבציה כללית}
היכולת לפתור בעיות קוונטיות באופן אנליטי היא מצומצמת ביותר. מערכות בודדות בלבד ניתנות ללכסון מדויק של ההמילטוניאן שלהן – כגון \LR{the Harmonic Oscillator} ואטום המימן.  

במערכות רבות בפיזיקה, המבנה מורכב מאוד, אך ניתן לזהות בהן הפרדה לרכיב "עיקרי" פתיר ועוד רכיב "קטן" המקלקל את הפתרון המלא. רכיב זה נקרא \textbf{הפרעה} (\LR{perturbation}).  

\begin{examplebox}{דוגמאות פיזיקליות להפרעות}
\begin{itemize}
\item \textbf{אטום מימן בשדה חשמלי (אפקט שטארק):} ההמילטוניאן העיקרי הוא אטום המימן, וההפרעה היא האינטראקציה עם השדה החשמלי $V = e\mathcal{E}z$.
\item \textbf{אטום מימן בשדה מגנטי (אפקט זימן):} ההפרעה היא האינטראקציה של המומנט המגנטי עם השדה $V = -\boldsymbol{\mu}\cdot\mathbf{B}$.
\item \textbf{צימוד ספין-מסלול:} תיקון רלטיביסטי קטן להמילטוניאן האטומי.
\item \textbf{אוסצילטור לא-הרמוני:} $V = \lambda x^3$ או $V = \lambda x^4$ מוסיפים על האוסצילטור ההרמוני.
\end{itemize}
\end{examplebox}

רעיון תורת ההפרעות הוא למנף את העובדה שהרכיב העיקרי פתיר, ולהוסיף תיקונים קטנים באופן שיטתי כסדרת חזקות בפרמטר הקטן.

\begin{definitionbox}{רעיון בסיסי — מבנה ההמילטוניאן}
נניח המילטוניאן מהצורה:
\[
\hat{H} = \hat{H}_0 + \lambda \hat{V}
\]
כאשר:
\begin{itemize}
\item $\hat{H}_0$ – ההמילטוניאן הלא-מופרע, פתיר לחלוטין עם ספקטרום ידוע $\{E_n^{(0)}, \ket{n^{(0)}}\}$.
\item $\hat{V}$ – אופרטור ההפרעה, בלתי תלוי בזמן. זהו אופרטור הרמיטי: $\hat{V}^\dagger = \hat{V}$.
\item $\lambda$ – פרמטר חסר-מימד המאפיין את עוצמת ההפרעה, כך ש־$\lambda \ll 1$.
\end{itemize}
\end{definitionbox}

\begin{notebox}{על הפרמטר $\lambda$}
הפרמטר $\lambda$ הוא כלי פורמלי שנועד לסדר את הפיתוח. בסוף החישוב לרוב נציב $\lambda = 1$ (או ששוקלים אותו בתוך $V$ מלכתחילה). הוא מאפשר לנו לעקוב אחר הסדר בסדרה: איברים עם $\lambda^1$ הם "סדר ראשון", איברים עם $\lambda^2$ הם "סדר שני", וכו'.
\end{notebox}

הגישה הפורמלית תתבסס על פיתוח סדרות טיילור בפונקציות הגל ובאנרגיות ביחס ל־$\lambda$.


\subsection{הנחות יסוד והקשר פיזיקלי}

לפני שנגזור את הנוסחאות, עלינו להבהיר את ההנחות שעליהן מבוססת התורה:

\begin{enumerate}
\item \textbf{ההמילטוניאן הלא-מופרע פתיר לחלוטין:}
\[
\hat{H}_0\ket{n^{(0)}} = E_n^{(0)}\ket{n^{(0)}}
\]
כאשר $\{\ket{n^{(0)}}\}$ מהווים בסיס אורתונורמלי מלא של מרחב הילברט:
\[
\braket{m^{(0)}|n^{(0)}} = \delta_{mn}, \qquad \sum_n \ket{n^{(0)}}\bra{n^{(0)}} = \hat{1}
\]

\item \textbf{אי-ניוון (\LR{Non-degeneracy}):} רמת האנרגיה שבה אנו מתמקדים אינה מנוונת:
\[
\forall m\neq n:\; E_m^{(0)} \neq E_n^{(0)}
\]
\textit{הערה:} אם קיים ניוון, יש להשתמש ב\textbf{תורת ההפרעות למצבים מנוונים} (ראו בהמשך).

\item \textbf{אנליטיות:} המערכות הן אנליטיות בפרמטר $\lambda$ — כלומר, הספקטרום $E_n(\lambda)$ ופונקציות הגל $\ket{n(\lambda)}$ הם פונקציות חלקות של $\lambda$, וסדרת הטיילור שלהם מתכנסת. זוהי הנחת עבודה שלא תמיד מתקיימת!
\end{enumerate}

\begin{notebox}{מתי ההנחות נכשלות?}
\textbf{מקרה 1 — רמות כמעט-מנוונות:} אם קיימות רמות אנרגיה עם הפרשים קטנים משמעותית מאלמנט המטריצה של ההפרעה:
\[
|E_n^{(0)} - E_k^{(0)}| \lesssim |\lambda V_{kn}|
\]
אז התורה נכשלת — המכנים הופכים קטנים והתיקונים "מתפוצצים".

\textbf{מקרה 2 — הפרעה "גדולה":} אם $\lambda$ אינו קטן באמת, או אם אלמנטי המטריצה של $V$ גדולים מדי.

\textbf{קריטריון תקפות:} יש לבדוק עבור כל $k \neq n$:
\[
\left|\lambda \frac{V_{kn}}{E_n^{(0)}-E_k^{(0)}}\right| \ll 1
\]
רק אם תנאי זה מתקיים, סדרת ההפרעות מתכנסת.
\end{notebox}


\section{פיתוח סדרות – מצבים ואנרגיות}

\subsection{הצבת האנזץ}
נרצה לפתור את משוואת הערכים העצמיים המלאה:
\[
\hat{H}\ket{n} = E_n\ket{n}
\]
כאשר $\hat{H} = \hat{H}_0 + \lambda\hat{V}$.

\textbf{הרעיון:} נפתח הן את המצב העצמי והן את האנרגיה כסדרות חזקות ב-$\lambda$:
\[
\ket{n} = \ket{n^{(0)}} + \lambda\ket{n^{(1)}} + \lambda^2\ket{n^{(2)}} + \lambda^3\ket{n^{(3)}} + \cdots
\]
\[
E_n = E_n^{(0)} + \lambda E_n^{(1)} + \lambda^2 E_n^{(2)} + \lambda^3 E_n^{(3)} + \cdots
\]

כאן:
\begin{itemize}
\item $\ket{n^{(0)}}$ ו-$E_n^{(0)}$ — המצב והאנרגיה הלא-מופרעים (ידועים).
\item $\ket{n^{(1)}}$, $E_n^{(1)}$ — תיקוני סדר ראשון (לא ידועים, נגזור).
\item $\ket{n^{(2)}}$, $E_n^{(2)}$ — תיקוני סדר שני (לא ידועים, נגזור).
\end{itemize}

\subsection{בחירת נרמול — שלב קריטי}

יש לנו חופש בבחירת הנורמליזציה של $\ket{n}$. נבחר נרמול \textbf{ביניים} נוח:
\[
\boxed{\braket{n^{(0)}|n} = 1}
\]

\textbf{משמעות:} הרכיב של $\ket{n}$ בכיוון $\ket{n^{(0)}}$ הוא בדיוק 1.

נפתח את התנאי הזה:
\begin{align}
\braket{n^{(0)}|n} &= \braket{n^{(0)}|n^{(0)}} + \lambda\braket{n^{(0)}|n^{(1)}} + \lambda^2\braket{n^{(0)}|n^{(2)}} + \cdots \\
&= 1 + \lambda\braket{n^{(0)}|n^{(1)}} + \lambda^2\braket{n^{(0)}|n^{(2)}} + \cdots = 1
\end{align}

כדי שזה יתקיים לכל $\lambda$ קטן, נדרש:
\[
\braket{n^{(0)}|n^{(k)}} = 0 \quad \text{לכל } k \geq 1
\]

\begin{notebox}{משמעות הנרמול}
בחירה זו איננה פיזיקלית אלא טכנית — היא מאפשרת הפרדה נקייה בין מקדמים בכל סדר של $\lambda$.

\textbf{שימו לב:} המצב $\ket{n}$ עם נרמול זה \textbf{אינו מנורמל ל-1}! הנורמה שלו היא:
\[
\braket{n|n} = 1 + O(\lambda^2)
\]
בסוף החישוב נצטרך לנרמל מחדש אם רוצים מצב פיזיקלי מנורמל.
\end{notebox}


\section{הכנסה למשוואת שרדינגר – גזירת תוצאות}

\subsection{הצבת הפיתוחים במשוואה}
נציב את פיתוחי הסדרות במשוואת הערכים העצמיים $\hat{H}\ket{n} = E_n\ket{n}$:
\[
(\hat{H}_0 + \lambda \hat{V})\Big(\ket{n^{(0)}} + \lambda\ket{n^{(1)}} + \lambda^2\ket{n^{(2)}} + \cdots\Big)
\]
\[
= \Big(E_n^{(0)} + \lambda E_n^{(1)} + \lambda^2 E_n^{(2)} + \cdots\Big)\Big(\ket{n^{(0)}} + \lambda\ket{n^{(1)}} + \lambda^2\ket{n^{(2)}} + \cdots\Big)
\]

נפתח את שני הצדדים ונאסוף איברים לפי חזקות $\lambda$:

\textbf{צד שמאל:}
\begin{align}
&\hat{H}_0\ket{n^{(0)}} + \lambda\Big(\hat{H}_0\ket{n^{(1)}} + \hat{V}\ket{n^{(0)}}\Big) \\
&+ \lambda^2\Big(\hat{H}_0\ket{n^{(2)}} + \hat{V}\ket{n^{(1)}}\Big) + O(\lambda^3)
\end{align}

\textbf{צד ימין:}
\begin{align}
&E_n^{(0)}\ket{n^{(0)}} + \lambda\Big(E_n^{(0)}\ket{n^{(1)}} + E_n^{(1)}\ket{n^{(0)}}\Big) \\
&+ \lambda^2\Big(E_n^{(0)}\ket{n^{(2)}} + E_n^{(1)}\ket{n^{(1)}} + E_n^{(2)}\ket{n^{(0)}}\Big) + O(\lambda^3)
\end{align}

כעת נשווה מקדמים בכל חזקה של $\lambda$.

\subsection{סדר אפס}
\[
\boxed{\hat{H}_0\ket{n^{(0)}} = E_n^{(0)}\ket{n^{(0)}}}
\]

זוהי פשוט משוואת הערכים העצמיים של ההמילטוניאן הלא-מופרע — לא חידוש, אלא אישור שהפיתוח עקבי.

\subsection{סדר ראשון}
\[
\hat{H}_0\ket{n^{(1)}} + \hat{V}\ket{n^{(0)}} = E_n^{(0)}\ket{n^{(1)}} + E_n^{(1)}\ket{n^{(0)}}
\]

נסדר מחדש:
\begin{equation}\label{eq:first_order}
\Big(\hat{H}_0 - E_n^{(0)}\Big)\ket{n^{(1)}} = E_n^{(1)}\ket{n^{(0)}} - \hat{V}\ket{n^{(0)}}
\end{equation}

\subsubsection{חילוץ תיקון האנרגיה מסדר ראשון}

נכפיל את משוואה \eqref{eq:first_order} משמאל ב-$\bra{n^{(0)}}$:
\[
\bra{n^{(0)}}\Big(\hat{H}_0 - E_n^{(0)}\Big)\ket{n^{(1)}} = E_n^{(1)}\underbrace{\braket{n^{(0)}|n^{(0)}}}_{=1} - \bra{n^{(0)}}\hat{V}\ket{n^{(0)}}
\]

\textbf{צד שמאל:}
\[
\bra{n^{(0)}}\hat{H}_0\ket{n^{(1)}} - E_n^{(0)}\braket{n^{(0)}|n^{(1)}}
\]

כעת נשתמש בעובדה ש-$\hat{H}_0$ הרמיטי, ולכן:
\[
\bra{n^{(0)}}\hat{H}_0 = \Big(\hat{H}_0\ket{n^{(0)}}\Big)^\dagger = \Big(E_n^{(0)}\ket{n^{(0)}}\Big)^\dagger = E_n^{(0)}\bra{n^{(0)}}
\]

(כי $E_n^{(0)}$ ממשי). לכן:
\[
\bra{n^{(0)}}\hat{H}_0\ket{n^{(1)}} = E_n^{(0)}\braket{n^{(0)}|n^{(1)}}
\]

ונקבל:
\[
E_n^{(0)}\braket{n^{(0)}|n^{(1)}} - E_n^{(0)}\braket{n^{(0)}|n^{(1)}} = 0
\]

\textbf{צד שמאל מתאפס!}

לכן:
\[
0 = E_n^{(1)} - \bra{n^{(0)}}\hat{V}\ket{n^{(0)}}
\]

\begin{resultbox}{תיקון מסדר ראשון לאנרגיה}
\[
\boxed{E_n^{(1)} = V_{nn} \equiv \bra{n^{(0)}}\hat{V}\ket{n^{(0)}}}
\]
\textbf{פרשנות פיזיקלית:} תיקון האנרגיה מסדר ראשון הוא פשוט \textbf{הערך התוחלתי של ההפרעה במצב הלא-מופרע}.
\end{resultbox}

\subsubsection{חילוץ תיקון המצב מסדר ראשון}

כעת נחלץ את $\ket{n^{(1)}}$. נכפיל את משוואה \eqref{eq:first_order} ב-$\bra{k^{(0)}}$ עבור $k \neq n$:
\[
\bra{k^{(0)}}\Big(\hat{H}_0 - E_n^{(0)}\Big)\ket{n^{(1)}} = E_n^{(1)}\underbrace{\braket{k^{(0)}|n^{(0)}}}_{=0} - \bra{k^{(0)}}\hat{V}\ket{n^{(0)}}
\]

\textbf{צד שמאל:}
\[
\bra{k^{(0)}}\hat{H}_0\ket{n^{(1)}} - E_n^{(0)}\braket{k^{(0)}|n^{(1)}}
\]

שוב, $\hat{H}_0$ הרמיטי ו-$\bra{k^{(0)}}\hat{H}_0 = E_k^{(0)}\bra{k^{(0)}}$:
\[
E_k^{(0)}\braket{k^{(0)}|n^{(1)}} - E_n^{(0)}\braket{k^{(0)}|n^{(1)}} = \Big(E_k^{(0)} - E_n^{(0)}\Big)\braket{k^{(0)}|n^{(1)}}
\]

לכן:
\[
\Big(E_k^{(0)} - E_n^{(0)}\Big)\braket{k^{(0)}|n^{(1)}} = -V_{kn}
\]

כאשר $V_{kn} \equiv \bra{k^{(0)}}\hat{V}\ket{n^{(0)}}$.

מכאן:
\[
\braket{k^{(0)}|n^{(1)}} = \frac{V_{kn}}{E_n^{(0)} - E_k^{(0)}} \quad (k \neq n)
\]

כעת נפרק את $\ket{n^{(1)}}$ בבסיס $\{\ket{m^{(0)}}\}$:
\[
\ket{n^{(1)}} = \sum_m \ket{m^{(0)}}\braket{m^{(0)}|n^{(1)}}
\]

נפריד את הסכום ל-$m=n$ ו-$m \neq n$:
\begin{itemize}
\item עבור $m = n$: $\braket{n^{(0)}|n^{(1)}} = 0$ (מתנאי הנרמול שבחרנו).
\item עבור $m \neq n$: $\braket{m^{(0)}|n^{(1)}} = \frac{V_{mn}}{E_n^{(0)} - E_m^{(0)}}$.
\end{itemize}

\begin{resultbox}{תיקון מסדר ראשון למצב העצמי}
\[
\boxed{\ket{n^{(1)}} = \sum_{k\neq n} \frac{V_{kn}}{E_n^{(0)}-E_k^{(0)}}\ket{k^{(0)}}}
\]
\textbf{פרשנות פיזיקלית:} ההפרעה "מערבבת" את המצב $\ket{n^{(0)}}$ עם מצבים אחרים $\ket{k^{(0)}}$. 
\begin{itemize}
\item ההשפעה גדולה יותר ככל שאלמנט המטריצה $V_{kn}$ גדול יותר.
\item ההשפעה גדולה יותר ככל שהמצבים קרובים באנרגיה (מכנה קטן).
\end{itemize}
\end{resultbox}

\begin{notebox}{מלכודת נפוצה — סימן המכנה}
שימו לב לסדר: המכנה הוא $E_n^{(0)} - E_k^{(0)}$ ולא להיפך. טעות בסימן היא שגיאה נפוצה מאוד!
\end{notebox}


\section{סדר שני – אנרגיות ונרמול}

\subsection{משוואת סדר שני}

נחזור לאיסוף האיברים מסדר $\lambda^2$:
\[
\hat{H}_0\ket{n^{(2)}} + \hat{V}\ket{n^{(1)}} = E_n^{(0)}\ket{n^{(2)}} + E_n^{(1)}\ket{n^{(1)}} + E_n^{(2)}\ket{n^{(0)}}
\]

נסדר:
\begin{equation}\label{eq:second_order}
\Big(\hat{H}_0 - E_n^{(0)}\Big)\ket{n^{(2)}} = E_n^{(1)}\ket{n^{(1)}} + E_n^{(2)}\ket{n^{(0)}} - \hat{V}\ket{n^{(1)}}
\end{equation}

\subsubsection{חילוץ תיקון האנרגיה מסדר שני}

נכפיל ב-$\bra{n^{(0)}}$:
\[
\underbrace{\bra{n^{(0)}}\Big(\hat{H}_0 - E_n^{(0)}\Big)\ket{n^{(2)}}}_{=0} = E_n^{(1)}\underbrace{\braket{n^{(0)}|n^{(1)}}}_{=0} + E_n^{(2)}\cdot 1 - \bra{n^{(0)}}\hat{V}\ket{n^{(1)}}
\]

(צד שמאל מתאפס מאותה סיבה כמו קודם, ו-$\braket{n^{(0)}|n^{(1)}}=0$ מתנאי הנרמול.)

לכן:
\[
E_n^{(2)} = \bra{n^{(0)}}\hat{V}\ket{n^{(1)}}
\]

נציב את הביטוי ל-$\ket{n^{(1)}}$:
\begin{align}
E_n^{(2)} &= \bra{n^{(0)}}\hat{V}\sum_{k\neq n}\frac{V_{kn}}{E_n^{(0)}-E_k^{(0)}}\ket{k^{(0)}} \\
&= \sum_{k\neq n}\frac{V_{kn}}{E_n^{(0)}-E_k^{(0)}}\bra{n^{(0)}}\hat{V}\ket{k^{(0)}} \\
&= \sum_{k\neq n}\frac{V_{nk}V_{kn}}{E_n^{(0)}-E_k^{(0)}}
\end{align}

מכיוון ש-$\hat{V}$ הרמיטי: $V_{nk} = V_{kn}^*$, ולכן $V_{nk}V_{kn} = |V_{kn}|^2$.

\begin{resultbox}{תיקון מסדר שני לאנרגיה}
\[
\boxed{E_n^{(2)} = \sum_{k\neq n}\frac{|V_{kn}|^2}{E_n^{(0)}-E_k^{(0)}}}
\]
\end{resultbox}

\begin{theorembox}{פרשנות פיזיקלית של $E_n^{(2)}$}
\begin{itemize}
\item \textbf{עבור מצב היסוד} ($n=0$): כל המכנים שליליים ($E_0^{(0)} < E_k^{(0)}$ לכל $k \neq 0$), ולכן:
\[
E_0^{(2)} < 0
\]
\textbf{תיקון סדר שני תמיד מוריד את אנרגית מצב היסוד!}

\item עבור מצבים מעוררים, הסימן תלוי באילו מצבים משפיעים יותר — אלה שמעליו או אלה שמתחתיו.

\item מצבים קרובים באנרגיה תורמים יותר (מכנה קטן).

\item מצבים עם אלמנט מטריצה גדול תורמים יותר (מונה גדול).
\end{itemize}
\end{theorembox}

\subsection{סיכום ביניים — האנרגיה עד סדר שני}

\begin{resultbox}{אנרגיה עד סדר שני}
\[
\boxed{E_n = E_n^{(0)} + \lambda V_{nn} + \lambda^2 \sum_{k\neq n}\frac{|V_{kn}|^2}{E_n^{(0)}-E_k^{(0)}} + O(\lambda^3)}
\]
\end{resultbox}

\subsection{נרמול מלא – גודל ההסטה הפיזיקלית}

המצב $\ket{n}$ שמצאנו אינו מנורמל (בחרנו $\braket{n^{(0)}|n}=1$ ולא $\braket{n|n}=1$). נחשב את הנורמה:
\begin{align}
\braket{n|n} &= \Big(\bra{n^{(0)}} + \lambda\bra{n^{(1)}} + \cdots\Big)\Big(\ket{n^{(0)}} + \lambda\ket{n^{(1)}} + \cdots\Big) \\
&= 1 + \lambda^2\braket{n^{(1)}|n^{(1)}} + O(\lambda^3)
\end{align}

(האיברים הלינאריים ב-$\lambda$ מתאפסים כי $\braket{n^{(0)}|n^{(1)}}=0$.)

נחשב את $\braket{n^{(1)}|n^{(1)}}$:
\begin{align}
\braket{n^{(1)}|n^{(1)}} &= \sum_{k\neq n}\sum_{m\neq n}\frac{V_{kn}^*V_{mn}}{(E_n^{(0)}-E_k^{(0)})(E_n^{(0)}-E_m^{(0)})}\braket{k^{(0)}|m^{(0)}} \\
&= \sum_{k\neq n}\frac{|V_{kn}|^2}{(E_n^{(0)}-E_k^{(0)})^2}
\end{align}

נגדיר מצב מנורמל:
\[
\ket{N} = Z_n^{1/2}\ket{n}, \qquad \braket{N|N} = 1
\]

כאשר:
\[
Z_n^{-1} = \braket{n|n} = 1 + \lambda^2\sum_{k\neq n}\frac{|V_{kn}|^2}{(E_n^{(0)}-E_k^{(0)})^2} + O(\lambda^3)
\]

\begin{theorembox}{קשר חשוב — המשקל הספקטרלי (\LR{Spectral Weight})}
\[
\boxed{Z_n = \frac{\partial E_n}{\partial E_n^{(0)}}}
\]
זהו קשר כללי ומעניין! הוא מראה שקבוע הנרמול $Z_n$ קשור לנגזרת האנרגיה לפי האנרגיה הלא-מופרעת.

\textbf{פרשנות:} $Z_n$ מודד את "מידת הישארות המצב המקורי" — כמה מהמצב $\ket{n}$ המלא נשאר בכיוון המקורי $\ket{n^{(0)}}$.
\end{theorembox}


\section{דוגמה מפורטת — אוסצילטור הרמוני עם הפרעה $x^2$}

נדגים את השיטה על דוגמה קונקרטית עם חישוב מלא.

\subsection{הגדרת הבעיה}
ניקח יחידות טבעיות $\hbar = m = \omega = 1$:
\[
\hat{H}_0 = -\frac{1}{2}\frac{d^2}{dx^2} + \frac{1}{2}x^2 = \frac{\hat{p}^2}{2} + \frac{\hat{x}^2}{2}
\]
\[
\hat{V} = x^2, \qquad \hat{H} = \hat{H}_0 + \lambda\hat{V}
\]

האנרגיות הלא-מופרעות:
\[
E_n^{(0)} = n + \frac{1}{2}, \qquad n = 0, 1, 2, \ldots
\]

\begin{notebox}{שימו לב — הבעיה הזו פתירה בדיוק!}
במקרה זה ניתן לפתור את הבעיה המלאה:
\[
\hat{H} = \frac{\hat{p}^2}{2} + \frac{1+2\lambda}{2}x^2 = \frac{\hat{p}^2}{2} + \frac{\omega'^2}{2}x^2
\]
כאשר $\omega' = \sqrt{1+2\lambda}$. האנרגיות המדויקות:
\[
E_n^{\text{exact}} = \left(n+\frac{1}{2}\right)\sqrt{1+2\lambda}
\]
נוכל להשוות לתוצאות תורת ההפרעות!
\end{notebox}

\subsection{חישוב אלמנטי המטריצה של $x^2$}

נזכיר את אופרטורי הסולם:
\[
\hat{a} = \frac{1}{\sqrt{2}}(\hat{x} + i\hat{p}), \qquad \hat{a}^\dagger = \frac{1}{\sqrt{2}}(\hat{x} - i\hat{p})
\]
\[
\hat{x} = \frac{1}{\sqrt{2}}(\hat{a} + \hat{a}^\dagger), \qquad \hat{p} = \frac{i}{\sqrt{2}}(\hat{a}^\dagger - \hat{a})
\]

לכן:
\begin{align}
\hat{x}^2 &= \frac{1}{2}(\hat{a} + \hat{a}^\dagger)^2 = \frac{1}{2}\Big(\hat{a}^2 + \hat{a}\hat{a}^\dagger + \hat{a}^\dagger\hat{a} + (\hat{a}^\dagger)^2\Big) \\
&= \frac{1}{2}\Big(\hat{a}^2 + (\hat{a}^\dagger)^2 + 2\hat{a}^\dagger\hat{a} + 1\Big)
\end{align}

כאשר השתמשנו ב-$[\hat{a},\hat{a}^\dagger]=1$ ולכן $\hat{a}\hat{a}^\dagger = \hat{a}^\dagger\hat{a} + 1$.

\textbf{כללי הבחירה:} מתוך המבנה של $\hat{x}^2$:
\begin{itemize}
\item $\hat{a}^2$ מוריד ב-2: $\bra{m}\hat{a}^2\ket{n} \neq 0$ רק אם $m = n-2$.
\item $(\hat{a}^\dagger)^2$ מעלה ב-2: $\bra{m}(\hat{a}^\dagger)^2\ket{n} \neq 0$ רק אם $m = n+2$.
\item $\hat{a}^\dagger\hat{a} + \frac{1}{2}$ אלכסוני.
\end{itemize}

\textbf{אלמנטי מטריצה ספציפיים:}

\underline{אלמנט אלכסוני:}
\[
\bra{n}x^2\ket{n} = \frac{1}{2}\bra{n}\Big(2\hat{a}^\dagger\hat{a} + 1\Big)\ket{n} = \frac{1}{2}(2n + 1) = n + \frac{1}{2}
\]

\underline{אלמנט $\bra{n+2}x^2\ket{n}$:}
\[
\bra{n+2}x^2\ket{n} = \frac{1}{2}\bra{n+2}(\hat{a}^\dagger)^2\ket{n} = \frac{1}{2}\sqrt{(n+1)(n+2)}
\]

\underline{אלמנט $\bra{n-2}x^2\ket{n}$:}
\[
\bra{n-2}x^2\ket{n} = \frac{1}{2}\bra{n-2}\hat{a}^2\ket{n} = \frac{1}{2}\sqrt{n(n-1)}
\]

\subsection{תיקון מסדר ראשון — מצב היסוד $n=0$}

\begin{resultbox}{תיקון מסדר ראשון לאנרגיית היסוד}
\[
E_0^{(1)} = \bra{0}x^2\ket{0} = 0 + \frac{1}{2} = \frac{1}{2}
\]
\end{resultbox}

\subsection{תיקון מסדר שני — מצב היסוד $n=0$}

עבור מצב היסוד, המצב היחיד שמתחבר אליו דרך $x^2$ הוא $\ket{2}$ (כי $\ket{-2}$ לא קיים):
\[
E_0^{(2)} = \sum_{k\neq 0}\frac{|V_{k0}|^2}{E_0^{(0)} - E_k^{(0)}} = \frac{|V_{20}|^2}{E_0^{(0)} - E_2^{(0)}}
\]

נחשב:
\[
V_{20} = \bra{2}x^2\ket{0} = \frac{1}{2}\sqrt{1\cdot 2} = \frac{\sqrt{2}}{2}
\]
\[
|V_{20}|^2 = \frac{1}{2}
\]
\[
E_0^{(0)} - E_2^{(0)} = \frac{1}{2} - \frac{5}{2} = -2
\]

\begin{resultbox}{תיקון מסדר שני לאנרגיית היסוד}
\[
E_0^{(2)} = \frac{1/2}{-2} = -\frac{1}{4}
\]
\end{resultbox}

\subsection{האנרגיה המקורבת ובדיקה}

האנרגיה עד סדר שני:
\[
E_0 \approx E_0^{(0)} + \lambda E_0^{(1)} + \lambda^2 E_0^{(2)} = \frac{1}{2} + \frac{\lambda}{2} - \frac{\lambda^2}{4}
\]

\textbf{בדיקה מול התוצאה המדויקת:}

נפתח את $E_0^{\text{exact}} = \frac{1}{2}\sqrt{1+2\lambda}$ בטיילור:
\[
\sqrt{1+2\lambda} = 1 + \lambda - \frac{\lambda^2}{2} + O(\lambda^3)
\]
\[
E_0^{\text{exact}} = \frac{1}{2}\Big(1 + \lambda - \frac{\lambda^2}{2} + O(\lambda^3)\Big) = \frac{1}{2} + \frac{\lambda}{2} - \frac{\lambda^2}{4} + O(\lambda^3)
\]

\begin{resultbox}{אימות התוצאה}
תורת ההפרעות נותנת תוצאה מדויקת עד לסדר $\lambda^2$! זו בדיקה חשובה של נכונות החישוב.
\end{resultbox}

\subsection{תרשים ויזואלי – מבנה רמות האנרגיה}

\begin{center}
\begin{tikzpicture}
\begin{axis}[
axis lines=middle,
ymin=0, ymax=5,
xmin=0, xmax=1.5,
ytick={0.5,1.5,2.5,3.5,4.5},
yticklabels={$E_0^{(0)}=\frac{1}{2}$,$E_1^{(0)}=\frac{3}{2}$,$E_2^{(0)}=\frac{5}{2}$,$E_3^{(0)}=\frac{7}{2}$,$E_4^{(0)}=\frac{9}{2}$},
xticklabels={},
xlabel={},
ylabel={אנרגיה},
title={רמות האוסצילטור ההרמוני}
]
\addplot[only marks, mark=-, mark size=10pt, thick] coordinates {(0.3,0.5)(0.3,1.5)(0.3,2.5)(0.3,3.5)(0.3,4.5)};
\node at (axis cs:0.3,0.1) {$H_0$};
\addplot[only marks, mark=-, mark size=10pt, thick, blue] coordinates {(1.0,0.6)(1.0,1.7)(1.0,2.8)(1.0,3.9)(1.0,5.0)};
\node at (axis cs:1.0,0.1) {$H_0+\lambda V$};
\draw[->, thick, red] (axis cs:0.45,0.5) -- (axis cs:0.85,0.6);
\end{axis}
\end{tikzpicture}
\end{center}

הרמות נעות כלפי מעלה (תיקון סדר ראשון חיובי) אך מצב היסוד "נדחף" למטה יחסית (תיקון סדר שני שלילי).


\section{אלגוריתם עבודה כללי — מתכון לפתרון בעיות}

\begin{enumerate}
\item \textbf{זהה את מבנה ההמילטוניאן:} כתוב $\hat{H} = \hat{H}_0 + \lambda\hat{V}$ כאשר $\hat{H}_0$ פתיר.

\item \textbf{פתור את הבעיה הלא-מופרעת:} מצא את כל האנרגיות $E_n^{(0)}$ והמצבים $\ket{n^{(0)}}$ של $\hat{H}_0$.

\item \textbf{בדוק ניוון:} האם יש רמות מנוונות? אם כן — עבור לתורת הפרעות למצבים מנוונים.

\item \textbf{חשב אלמנטי מטריצה:} חשב $V_{kn} = \bra{k^{(0)}}\hat{V}\ket{n^{(0)}}$ לכל הזוגות הרלוונטיים.

\textit{טיפ:} חפש כללי בחירה! לעיתים קרובות סימטריות גורמות לאלמנטים רבים להתאפס.

\item \textbf{חשב תיקון סדר ראשון לאנרגיה:}
\[
E_n^{(1)} = V_{nn}
\]

\item \textbf{חשב תיקון סדר ראשון למצב:}
\[
\ket{n^{(1)}} = \sum_{k\neq n}\frac{V_{kn}}{E_n^{(0)}-E_k^{(0)}}\ket{k^{(0)}}
\]

\item \textbf{חשב תיקון סדר שני לאנרגיה (אם נדרש):}
\[
E_n^{(2)} = \sum_{k\neq n}\frac{|V_{kn}|^2}{E_n^{(0)}-E_k^{(0)}}
\]

\item \textbf{בדוק תקפות:} וודא שלכל $k \neq n$:
\[
\left|\lambda\frac{V_{kn}}{E_n^{(0)} - E_k^{(0)}}\right| \ll 1
\]
אם התנאי לא מתקיים, התוצאות אינן אמינות.

\item \textbf{נרמל אם צריך:} אם דרוש מצב מנורמל פיזיקלית, חשב $Z_n$ ונרמל.
\end{enumerate}

\begin{notebox}{טיפים מעשיים}
\begin{itemize}
\item \textbf{כללי בחירה:} נצל סימטריות! אם $\hat{V}$ זוגי ב-$x$, הוא מחבר רק מצבים עם אותה זוגיות. אם $\hat{V}$ אי-זוגי, הוא מחבר מצבים עם זוגיות הפוכה.
\item \textbf{סכומים אינסופיים:} לעיתים הסכום מתכנס מהר ומספיק לקחת מספר איברים ראשונים.
\item \textbf{בדיקת יחידות:} ודא שליחידות יש משמעות — $E^{(1)}$ צריך להיות באותן יחידות כמו $E^{(0)}$.
\end{itemize}
\end{notebox}


\section{מתי התורה נכשלת ומה עושים אז?}

\subsection{מקרי כישלון}

\begin{notebox}{כישלון התורה — סימנים ופתרונות}
התורה נכשלת או נותנת תוצאות לא אמינות במקרים הבאים:

\textbf{1. רמות מנוונות או כמעט-מנוונות:}
\begin{itemize}
\item אם $E_m^{(0)} = E_n^{(0)}$ עבור $m \neq n$, המכנים מתאפסים והנוסחאות מתפוצצות.
\item אם $|E_m^{(0)} - E_n^{(0)}| \ll |\lambda V_{mn}|$, התיקונים גדולים מדי והסדרה לא מתכנסת.
\item \textbf{פתרון:} תורת הפרעות למצבים מנוונים (\LR{Degenerate Perturbation Theory}).
\end{itemize}

\textbf{2. הפרעה "גדולה":}
\begin{itemize}
\item אם $\lambda$ אינו קטן באמת, או אם אלמנטי המטריצה גדולים.
\item \textbf{פתרון:} שיטות לא-פרטורבטיביות, לכסון מספרי, שיטות וריאציוניות.
\end{itemize}

\textbf{3. התכנסות איטית או אי-התכנסות:}
\begin{itemize}
\item לעיתים הסדרה מתכנסת אך לאט מאוד, ונדרשים סדרים גבוהים רבים.
\item לעיתים הסדרה אסימפטוטית ולא מתכנסת כלל (למשל, $V = \lambda x^4$ באוסצילטור הרמוני).
\item \textbf{פתרון:} טכניקות ריסום (\LR{Resummation}), רנורמליזציה.
\end{itemize}
\end{notebox}

\subsection{תורת הפרעות למצבים מנוונים — רעיון בסיסי}

כאשר יש ניוון, נניח $g$-פולד (כלומר, $g$ מצבים עם אותה אנרגיה $E^{(0)}$), הנוסחאות הרגילות נכשלות כי המכנה מתאפס.

\textbf{הרעיון:} במקום לטפל בכל מצב בנפרד:
\begin{enumerate}
\item הגדר תת-מרחב מנוון: $\mathcal{D} = \text{span}\{\ket{1^{(0)}}, \ket{2^{(0)}}, \ldots, \ket{g^{(0)}}\}$
\item לכסן את $\hat{V}$ \textbf{בתוך} תת-המרחב הזה.
\item הערכים העצמיים של המטריצה $g \times g$:
\[
V_{ij} = \bra{i^{(0)}}\hat{V}\ket{j^{(0)}}, \quad i,j = 1,\ldots,g
\]
נותנים את תיקוני האנרגיה מסדר ראשון.
\item הוקטורים העצמיים קובעים את הקומבינציות הלינאריות "הנכונות" של מצבי הבסיס.
\end{enumerate}

\textit{נושא זה יורחב בהרצאה הבאה.}

\subsection{דוגמה לכישלון — אטום מימן בשדה חשמלי (אפקט שטארק)}

באטום המימן, הרמה $n=2$ מנוונת ארבע-פולד (מצבים $2s, 2p_x, 2p_y, 2p_z$ — כולם עם אנרגיה $E_2 = -\frac{1}{8}\text{Ry}$).

ההפרעה $\hat{V} = e\mathcal{E}z$ מחברת בין מצבים אלה. לא ניתן להשתמש בנוסחה הרגילה, ויש ללכסן את המטריצה $4 \times 4$ של $\hat{V}$ בתת-מרחב הזה.


\section{סיכום עיקרי הנקודות}

\begin{resultbox}{תמצית תורת ההפרעות הלא-תלויה בזמן}
\textbf{המבנה:}
\[
\hat{H} = \hat{H}_0 + \lambda\hat{V}, \qquad \lambda \ll 1
\]

\textbf{תיקון סדר ראשון לאנרגיה:}
\[
E_n^{(1)} = \bra{n^{(0)}}\hat{V}\ket{n^{(0)}} = V_{nn}
\]

\textbf{תיקון סדר ראשון למצב:}
\[
\ket{n^{(1)}} = \sum_{k\neq n}\frac{V_{kn}}{E_n^{(0)}-E_k^{(0)}}\ket{k^{(0)}}
\]

\textbf{תיקון סדר שני לאנרגיה:}
\[
E_n^{(2)} = \sum_{k\neq n}\frac{|V_{kn}|^2}{E_n^{(0)}-E_k^{(0)}}
\]

\textbf{תנאי תקפות:}
\[
\left|\lambda\frac{V_{kn}}{E_n^{(0)}-E_k^{(0)}}\right| \ll 1 \quad \forall k \neq n
\]
\end{resultbox}

\subsection*{בדיקות הבנה}

\begin{enumerate}
\item מדוע תיקון סדר שני לאנרגיית מצב היסוד תמיד שלילי?

\item עבור אוסצילטור הרמוני עם הפרעה $V = \lambda x^3$:
\begin{itemize}
\item מה תיקון סדר ראשון לאנרגיה? (רמז: חשבו על זוגיות)
\item אילו מצבים מתחברים דרך $x^3$?
\end{itemize}

\item תנו דוגמה למערכת בה תורת ההפרעות אינה מתאימה, והסבירו מדוע.

\item הוכיחו שאם $\hat{V}$ קומוטטיבי עם $\hat{H}_0$, אז $\ket{n^{(1)}} = 0$ ותיקוני האנרגיה הם בדיוק הערכים העצמיים של $\hat{V}$.
\end{enumerate}

